\begin{table}
\centering
\footnotesize
\begin{tabular}{@{}l l c c@{}}
\toprule
\textbf{Branche} & \textbf{Modèle} & \textbf{F1-macro} & \textbf{Accuracy} \\
\midrule
A — Stats       & \ModelAStats   & \FOneAStats   & \AccAStats \\
B — Tex(opt)    & \ModelBTex    & \FOneBTex    & \AccBTex \\
C — DWT(opt)    & \ModelCDwt    & \FOneCDwt    & \AccCDwt \\
D — HOG         & \ModelDHog    & \FOneDHog    & \AccDHog \\
A+B             & \ModelApB     & \FOneApB     & \AccApB \\
B+C(opt)        & \ModelBpC     & \FOneBpC     & \AccBpC \\
A+B+C(opt)      & \ModelApBpC   & \FOneApBpC   & \AccApBpC \\
\textbf{Full(opt)} & \textbf{\ModelFull} & \textbf{\FOneFull} & \textbf{\AccFull} \\
\bottomrule
\end{tabular}
\caption{Meilleure combinaison modèle/branche sur le split test (hold-out).}
\end{table}
